\section{Boundary attainment and finite optimality witnesses}\label{sec:primal-dual}
An upper construction and a lower inequality do not by themselves characterize which strategies are optimal. We now close the constrained risk frontier at its feasible boundary, derive its exact variational dual at a fixed dictionary, and identify finite witnesses for each compact-model prefix problem. Cost and task witnesses are allowed to be different physical preparations in different models. Their common-program Bellman certificate is necessary and sufficient at the finite level, and, together with the preceding paid completion, certifies arbitrarily accurate strategies at every interior budget.

The convex-analytic ingredients are classical: finite-support semi-infinite duality is developed in \cite{Charnes63,Shapiro09}, and compactness of strategic and occupation measures has an extensive control-theoretic history \cite{Balder89,Dufour25}. We prove the needed statements in this experiment's coordinates. We do not claim a new general separation or Carath\'eodory theorem. The distinction between a fixed dictionary and a hard bound on the number of freely chosen decoder values is essential: the latter can destroy convexity even in physical dimension one.

\subsection{The exact frontier, including its left endpoint}
For this subsection only, the model family need not be compact, time homogeneous, or finite-dimensional. It suffices that all models share a countable, finitely branching action--report tree with nonempty finite legal action sets at open nodes. Costs have finite conditional means at each finite history, are nonnegative, and have the same positive conditional mean floor $c_*$. Terminal losses lie in $[0,L]$. Initial preparations may be absorbed into the model index. One common observer must commit exactly once and satisfy its expected-cost bound in every model. The decision set $D$ is compact metric and the terminal losses have a common decision modulus. These conditions are satisfied by the primitive interfaces above; they specify the wider scope used here without assuming a free posterior or a common physical reset.

Use the common sequence weights of Lemma~\ref{lem:common-plan}. A fixed ordered decoder tuple $C$ may have repeated entries. Let $\mathcal P_N(C)$ denote the common infinite plans of uniform expected cost at most $N$. Let $q_{C,T}(N)$ be the infimum of worst truncated terminal loss over common depth-$T$ plans with the one-call reserve $C_{\theta,T}+c_*S_{\theta,T}\le N$ in every model and preparation. For an operator interface this inequality means the corresponding Loewner inequality. Define $q_{m,T}$ by optimizing the decoder tuple in $D^m$, and likewise define $R_m$. A finite infeasible problem has value $+\infty$.

\begin{theorem}[Closed constrained frontier and boundary convergence]\label{thm:boundary-attainment}
Every nonempty $\mathcal P_N(C)$ is compact in common-plan coordinates. Terminal-risk functions converge uniformly to their depth-$T$ truncations on this set, with error at most
\[
 L\min\{1,(N/c_*-1)/T\}\qquad(T\ge1).
\]
Thus the optimal risk is attained at every feasible finite budget, including the least feasible budget when one exists. The same assertion holds with a compact decoder tuple of prescribed width $m$.

A budget is feasible if and only if all its finite reserved problems are feasible. At every feasible $N$,
\begin{equation}\label{eq:boundary-convergence}
 q_{C,T}(N)\uparrow R_C(N),\qquad q_{m,T}(N)\uparrow R_m(N).
\end{equation}
Both risk functions are nonincreasing and right-continuous on their feasible budget ray. For a fixed dictionary, $R_C$ is convex and hence continuous at every interior feasible budget. No convexity of $R_m$ is asserted. Boundary convergence gives neither a computable rate nor a finite-memory optimum at the least budget.
\end{theorem}
\begin{proof}
The infinite common-plan space is a closed subset of the countable product of the intervals $[0,1]$, specified by the common flow equalities. It is compact and metrizable. Every prefix cost, survival mass and truncated terminal loss is continuous in its finitely many sequence coordinates for each model. Intersecting all prefix reserve constraints over every model and every depth gives a closed set.

On that intersection, the cost floor and monotonicity of survival give
$C_T+c_*S_T\ge c_*[1+T S_T]$. Hence the displayed uniform survival bound holds, and survival tends to zero. Nonnegative prefix costs increase to the actual cost; their bounds imply a proper common plan of cost at most $N$. Conversely a proper cost-$N$ plan satisfies every reserve because each open path owes a further paid call. This identifies the closed intersection with $\mathcal P_N(C)$, without assuming continuity of total cost or a uniform expected remaining-cost tail.

The nonnegative omitted terminal loss is at most $L$ times survival. A finite prefix loss is uniformly continuous in its finite plan coordinates, uniformly over models: its coefficients are bounded by $L$, and there are finitely many terminal paths. For a decoder tuple, the common decision modulus supplies the additional uniform continuity. Thus the full worst-model risk is a uniform limit of continuous functions on the compact cost-bounded set. It attains its minimum. The product with $D^m$ is compact as well; output indices, not unknown model-dependent values, are the terminal choices.

A later feasible prefix restricts to an earlier one while preserving its reserve and decreasing its accumulated nonnegative loss. Thus the finite values are nondecreasing and bounded above by the actual value. Conversely choose optimizers at depths tending to infinity, extend them by arbitrary legal rows, and take a coordinatewise convergent subsequence, also in $D^m$ when needed. At each fixed earlier depth the limiting reserve and terminal-loss constraints hold in every model. The limit is proper by the preceding argument, and its terminal risk is at most the increasing limit of the finite values. This proves \eqref{eq:boundary-convergence}. The same compact finite-intersection argument without the loss constraints proves the feasibility equivalence. Choosing a sequence of feasible budgets decreasing to their infimum also proves attainment of the least finite feasible cost.

For right-continuity, use optimizers at budgets $N_j\downarrow N$, all bounded by $N+1$, and take a convergent subsequence in that compact cost-bounded set. Its limit is feasible at $N$, and uniform terminal-tail control identifies the limiting risk. Monotonicity supplies the reverse inequality. Finally, two common plans with the same dictionary may be mixed in sequence coordinates. Their costs and losses mix linearly before taking worst models, so the worst cost and risk mix sublinearly. This proves convexity of $R_C$. Such mixing is a common behavioral policy by Lemma~\ref{lem:common-plan}; a finite implementation must count its live branch roles as in the supplement. Mixing two optimized width-$m$ dictionaries can use $2m$ values, so this argument does not prove convexity of $R_m$.
\end{proof}

For a fixed dictionary, write $\mathfrak C(P)$ and $\mathfrak V(P)$ for a plan's worst expected cost and worst terminal risk, and let $\mathcal P_{\rm fin}$ contain the proper common plans of finite uniform cost. The following is a scalar resource dual, not an interchange of a physical model and an observer.
\begin{corollary}[Exact resource-price dual]\label{cor:resource-price}
If some finite budget is feasible, then for every $N\ge N_{\min}$,
\begin{equation}\label{eq:resource-price}
 R_C(N)=\sup_{\lambda\ge0}\inf_{P\in\mathcal P_{\rm fin}}
       \{\mathfrak V(P)+\lambda[\mathfrak C(P)-N]\}.
\end{equation}
At $N>N_{\min}$ the supremum is attained. If a common plan of cost at most $\kappa<N$ is known, an attaining multiplier satisfies $0\le\lambda\le L/(N-\kappa)$. Every primal optimum then minimizes the penalized expression and obeys
$\lambda[\mathfrak C(P)-N]=0$. At $N=N_{\min}$ the supremum remains exact but need not be attained by a finite multiplier. These statements concern one fixed dictionary, not its nonconvex width-constrained union.
\end{corollary}
\begin{proof}
Weak duality follows by testing any feasible plan. On the interior the finite convex nonincreasing function $R_C$ has a supporting line of slope $-\lambda\le0$ at $N$. For every finite-cost plan of cost $n$, its risk is at least $R_C(n)$, which is at least $R_C(N)-\lambda(n-N)$. A primal optimum gives the reverse bound and complementary slackness. Applying the supporting inequality at $\kappa$ gives $\lambda(N-\kappa)\le R_C(\kappa)-R_C(N)\le L$. Supporting lines exist because the left and right secant slopes of a finite convex function are ordered and locally bounded. At the left endpoint choose $N_j\downarrow N_{\min}$ from the interior. The attaining interior multiplier gives a dual value at $N_{\min}$ at least $R_C(N_j)$. Right-continuity and weak duality prove endpoint equality without a bounded-multiplier assumption.
\end{proof}


\subsection{A complete decision comparison of performed stopped laws}
Fix the actual decoder dictionary $C$, rather than only its cardinality. Let $\Omega$ be the measurable disjoint union of terminal action--report words, committed values in $C$, and a prospective audit outcome in a common standard Borel space. The audit kernel after each such word and commitment is fixed in each model. Write $\mathsf P_\theta^P$ for the complete performed stopped law of a legal plan $P$. It includes the stopping call and the raw visible word, not just the final audit marginal. For a specified target family of probability laws $\mathsf T_\theta$ on this same space define
\[
 \delta_{N,C}(\mathsf T)
 =\inf_{P\in\mathcal P_N(C)}\sup_\theta
                \TV(\mathsf P_\theta^P,\mathsf T_\theta),
 \qquad \TV(\mu,\nu)=\sup_A|\mu(A)-\nu(A)|.
\]
The target is data, not an extra executable oracle. Its feasibility in the source interface is the question being tested.

\begin{theorem}[Paid complete-score randomization criterion]\label{thm:complete-stopped-order}
Under the common finite-branching and positive-cost hypotheses of Theorem~\ref{thm:boundary-attainment}, at any feasible finite $N$, the set of families $\{\mathsf P_\theta^P\}_\theta$ is compact and convex in the uniform total-variation metric. Hence the infimum defining $\delta_{N,C}$ is attained, and
\begin{equation}\label{eq:complete-stopped-dual}
 \delta_{N,C}(\mathsf T)
 =\sup_{\pi,g}\left[
 \min_{P\in\mathcal P_N(C)}\sum_{\theta\in\operatorname{supp}\pi}
           \pi_\theta\E_{\mathsf P_\theta^P}g_\theta
 -\sum_{\theta\in\operatorname{supp}\pi}
           \pi_\theta\E_{\mathsf T_\theta}g_\theta\right].
\end{equation}
Here $\pi$ ranges over finitely supported model probabilities and each $g_\theta:\Omega\to[0,1]$ is a finite-valued measurable score. The inner policies remain budget-feasible in \emph{every} model, not only the support of $\pi$. In particular exact common realization of the target at the prescribed budget and dictionary is equivalent to its dominating the optimal source score in every such finite decision test.

On the compact primitive interface, let one common finite-label fallback have conditional cost at most $\kappa<N$. For every $T\ge1$ and every $P\in\mathcal P_N(C)$ there is a common finite-label $P_T\in\mathcal P_N(C)$ satisfying
\begin{equation}\label{eq:stopped-law-density}
 \sup_\theta\TV(\mathsf P_\theta^{P_T},\mathsf P_\theta^P)
 \le D_T+\frac{e_T}{N-\kappa+e_T},\qquad
 D_T=\min\{1,(N/c_*-1)/T\},\quad e_T=(\kappa-c_*)D_T.
\end{equation}
The prior finite-prefix and fallback live-label bounds apply. Therefore finite-label laws are dense at the same budget at every strictly interior feasible budget, and finite-score lower witnesses and actual finite-label realizations approximate the exact deficiency. This is a comparison of the declared performed laws, not of every quantum channel with arbitrary ancillary experiments.
\end{theorem}
\begin{proof}
At depth at most $T$ there are finitely many terminal action--report words and output indices. Their contributions to a stopped law are fixed physical audit measures multiplied by common sequence weights. For two plans, the total-variation difference of these finite contributions is at most the sum of the absolute differences of those weights, uniformly over models, since each physical measure has mass at most one. The two omitted positive tails have mass at most $D_T$ each on the budget set. First choose $T$ large and then make the finitely many coordinates close. This proves continuity of the law-family map into uniform total variation. Compactness of $\mathcal P_N(C)$ gives compactness of its image; linearity in sequence weights gives convexity. Distance to any fixed target family is a continuous function in this metric, so its minimum is attained. This argument does not need continuity of total expected cost or uniform integrability of its tail.

For equal-mass probability measures, $\TV(\mu,\nu)=\sup_A[\mu(A)-\nu(A)]$, since complements reverse the sign. Thus each test functional
$F_{\theta,A}(P)=\mathsf P_\theta^P(A)-\mathsf T_\theta(A)$ is continuous and affine on the compact convex plan set. The right side of \eqref{eq:complete-stopped-dual} is at most the deficiency by the elementary minimax inequality. If $r<\delta_{N,C}$, the closed sets $\{P:F_{\theta,A}(P)\le r\}$ have empty total intersection. Compactness gives a finite subcollection with empty intersection. Its attainable test vectors form a compact convex set disjoint from the lower orthant $(-\infty,r]^q$. Strict separation gives nonnegative weights summing to one, with the weighted test value greater than $r$ for every feasible plan. Nonnegativity follows because a negative coefficient would make the supremum on that lower orthant infinite. Grouping tests with the same model yields a finitely supported $\pi$ and a convex combination of finitely many indicators as $g_\theta$. The latter is finite-valued. This proves the reverse inequality by letting $r$ increase to the deficiency. Zero deficiency implies exact realization by attainment, not only approximation.

For \eqref{eq:stopped-law-density}, keep $P$ unchanged through the depth-$T$ commitment decision, and execute the real fallback on its open paths. Coupling the original and completed experiments with the same calls and current coins until then makes their complete stopped laws identical unless the original path is still open. The total-variation discrepancy is therefore at most $D_T$, including all raw history and audit coordinates. The completed cost is at most $N+e_T$, because its reserved prefix already paid $c_*$ for each open path. Mix this completion with the common fallback with probability $p=(N-\kappa)/(N-\kappa+e_T)$. The expected cost is at most $N$ in every model. Mixing changes a probability law by at most $1-p$ in total variation. The triangle inequality proves the displayed bound. Its finite implementation counts the actual finite tree, the fallback and the initial branch role; the same dictionary is used throughout. The compact-family cost theorem provides such a fallback at every $N>N_{\min}$. Applying this approximation to an attaining deficiency minimizer proves the final assertion.
\end{proof}

The decision-separation principle belongs to the Blackwell--Le Cam tradition \cite{Blackwell,Jencova16}. The additional work here is the compactness of the \emph{performed} law family under a uniform paid stopping constraint, together with finite-label approximation without increasing that constraint. It requires all bounded measurable scores of that performed outcome; localized affine-square tests alone do not suffice. The conclusion does not grant a free stateful simulator, an uncharged audit on rejected calls, or control of unbounded expected duration by total variation. The latter is controlled by the separate budget throughout. If the dictionary itself varies, laws recording different numerical readouts need not vary continuously in total variation; the general-loss width theorem and the fixed-dictionary comparison are distinct statements.

\subsection{Sparse spectral certificates on compact uncertainty}
Return to the continuous compact finite-dimensional interface of Section~\ref{sec:compact-risk}. Fix a dictionary and horizon. Let $\mathcal X_T$ be the common finite sequence-plan polytope and let $k_T$ be its affine dimension, after eliminating flow equalities. The same statement applies to a common tree with compulsory terminal nodes, or to a prescribed face of its plan polytope; $k_T$ then denotes the affine dimension of that polytope. For $z=(\theta,u)$ with $\theta\in\Theta$ and $\|u\|=1$, put
\[
 c_z(x)=\langle u,(C_{\theta,T}(x)+c_*S_{\theta,T}(x))u\rangle,
 \qquad v_z(x)=\langle u,V_{\theta,T}(x)u\rangle.
\]
These are continuous families of affine functions of $x$; $0\le v_z(x)\le L$. Testing all such $z$ is equivalent to the operator constraints. Write $q_T(N)$ for the resulting compact-model finite optimum.

\begin{theorem}[Exact finite witness and complementary slackness]\label{thm:sparse-witness}
Suppose some $x_0\in\mathcal X_T$ satisfies $c_z(x_0)\le\kappa<N$ for every $z$. Then there are risk atoms $z_j$, cost atoms $w_i$, and nonnegative weights $\beta_j,\alpha_i$ such that
\begin{align}
 &\sum_j\beta_j=1,\qquad
 \lambda:=\sum_i\alpha_i\le\frac{L}{N-\kappa},\qquad
 \#\{j:\beta_j>0\}+\#\{i:\alpha_i>0\}\le k_T+1,\label{eq:sparse-size}\\
 &q_T(N)=\min_{x\in\mathcal X_T}
 \left\{\sum_j\beta_jv_{z_j}(x)+
                    \sum_i\alpha_i[c_{w_i}(x)-N]\right\}.\label{eq:sparse-dual}
\end{align}
For any fixed optimal plan $x_*$ these atoms can be chosen active there:
\[
 \beta_j>0\Rightarrow v_{z_j}(x_*)=q_T(N),\qquad
 \alpha_i>0\Rightarrow c_{w_i}(x_*)=N.
\]
Their preparations are pure; active vectors are respectively eigenvectors of the loss operator at $q_T$ and the reserved-cost operator at $N$. Conversely, a feasible plan, such active atoms and minimization of the displayed affine expression certify optimality. The total atom bound is sharp for each affine dimension $k_T$. It concerns dual witnesses, not the number of controller labels or a fixed horizon-independent adversarial prior.
\end{theorem}
\begin{proof}
Set $v(x)=\max_zv_z(x)$ and $c(x)=\max_zc_z(x)$. Compactness gives attained maxima. The finite value as a function of its budget is convex and nonincreasing. A supporting line at $N$, which is interior because $\kappa<N$, yields a multiplier $\lambda\ge0$ for which $x_*$ minimizes $v+\lambda(c-N)$ and $\lambda(c(x_*)-N)=0$, exactly as in Corollary~\ref{cor:resource-price}.

For completeness, the finite polytope argument underlying its stationarity is proved here. The directional derivative of a compact maximum of affine functions is the maximum of their active gradients: take a maximizing index along a sequence of positive steps, pass to a convergent subsequence, and compare with every active index. Let $A$ be the compact convex sum of the active risk-gradient hull and $\lambda$ times the active cost-gradient hull; omit the second term if $\lambda=0$. If $A$ did not intersect the negative normal cone of $\mathcal X_T$ at $x_*$, separation of a compact convex set and a closed convex cone would give a tangent direction $h$ with $\max_{a\in A}a\cdot h<0$. For a polytope, every vector in this tangent cone gives a feasible sufficiently short step: impose its finitely many active linear inequalities, then choose the step short enough for the inactive ones. The directional derivative would be negative, a contradiction. Thus there are active finite convex combinations whose sum $g$ satisfies $g\cdot(x-x_*)\ge0$ for all $x\in\mathcal X_T$. Convex combinations in a finite-dimensional compact hull can be chosen finite by eliminating affine dependencies.

In the $k_T$-dimensional affine coordinates, append one coordinate to the active gradients. Represent $(g,1)$ as a conic sum of $(\nabla v_z,1)$ for risk atoms and $(\nabla c_w,0)$ for cost atoms. If more than $k_T+1$ generators have positive coefficients, a linear dependence permits subtracting a suitable multiple of that dependence until one coefficient vanishes and none becomes negative. Repeat. This preserves $g$ and risk mass one, and all remaining atoms stay active. Their affine combination is minimized at $x_*$. Its value there is $q_T$, proving \eqref{eq:sparse-dual}. Evaluating that minimum at $x_0$ gives $q_T\le L-\lambda(N-\kappa)$, hence the bound in \eqref{eq:sparse-size}, even if compression changed the original multiplier.

For sufficiency, at any feasible $x$ the expression in braces is at most $v(x)$, while at the specified active minimizing plan it equals its risk. This sandwiches the optimum. If a positive semidefinite matrix $NI-C$ has zero quadratic form at a unit vector, it annihilates that vector; the spectral claims follow, and similarly for risk. Finally take a compulsory one-call scalar experiment with $k+1$ actions and $k+1$ models, unit cost, one decoder and loss $1$ exactly when action and model indices coincide. With commitment required after that call, its plan polytope is the $k$-simplex and its risk is $\max_i x_i$. At any $N>1$, cost atoms cannot be active. The value $1/(k+1)$ requires the risk weights all equal $1/(k+1)$, since the dual minimum is their smallest weight. Every one of the $k+1$ atoms is necessary. This proves sharpness within physical experiments, not just an abstract cone.
\end{proof}

Finite-support duality and the conic reduction are classical \cite[Theorem 4.4]{Shapiro09}. The operational content is their use with the actual common plan, the open-path reserve, distinct spectral cost and score preparations, and a completion at the same resource budget. An arbitrary coarse grid is not a replacement for this certificate: active models may lie in the interior of a parameter set, as Example~\ref{ex:vertex-risk} shows.

\begin{proposition}[Common Bellman evaluation and equality of used rows]\label{prop:dual-bellman}
For any finite atom lists with risk mass one, the right side of \eqref{eq:sparse-dual} is evaluated exactly by a backward recursion on the single action--report tree. At each productive report it minimizes between a decision in the fixed dictionary and continuation; nonproductive reports must continue. The open depth-$T$ cost is the actual one-call reserve. There is a deterministic finite plan minimizing this dual expression, but it need not satisfy the original robust budget and risk constraints. A randomized plan minimizes it if and only if every choice of positive sequence weight has zero backward slack at its reached node.
\end{proposition}
\begin{proof}
Propagate, for each atom, its unnormalized physical state along each raw action--report word, without multiplying by controller probabilities. At a pre-call word $h$, the immediate cost for action $a$ is $\sum_i\alpha_i\tr(G_{w_i,a}\rho_{w_i,h})$. At its productive report $y$, stopping with $r$ contributes $\sum_j\beta_j\tr(F_{z_j,a,y,r}\rho_{z_j,hay})$. Continuation contributes the next node's recursively computed value. At an open leaf of depth $T$, assign $c_*\sum_i\alpha_i\tr\rho_{w_i,h}$. Add over reports and minimize over the common action. Subtract $\lambda N$ once at the root. This evaluates the desired affine functional by backward induction; the full formula is recorded in the supplement.

Each minimization admits a deterministic choice because its alternatives are finite. Keeping the chosen current command or report never reveals the model: the propagated matrices are offline analytic coefficients for a common tree. At every row, its randomized value minus its minimum is a nonnegative slack. Backward substitution expresses the root excess as the sum of those slacks times the corresponding nonnegative sequence weights. It vanishes precisely when every positively used alternative has zero slack. A dual-minimizing plan alone may violate the individual model constraints; Theorem~\ref{thm:sparse-witness} requires their separate primal verification.
\end{proof}

\begin{corollary}[Finite certificates for the full stopped optimum]\label{cor:complete-certificate}
Fix $N>N_{\min}$ and a dictionary $C$ on the compact primitive interface. For every $\eta>0$ there is a finite horizon, at most $k_T+1$ spectral witness atoms, and a common proper finite-label observer at cost at most $N$ such that the Bellman-evaluated dual lower and that observer's worst risk differ by at most $\eta$. The same actual dictionary is used at every history and by the fallback. With compact optimized output width, attainment still holds, but this dual is conditional on the dictionary; global certification requires the separate dictionary optimization or a certified decision net. No minimax swap over dictionaries is authorized.
\end{corollary}
\begin{proof}
The compact-family cost hierarchy supplies a finite fallback of uniform cost $\kappa<N$. Its prefix is strictly feasible for every reserved finite problem. Apply Theorem~\ref{thm:sparse-witness}. The compact-model problem imposes all model constraints, so completion and initial mixing introduce no model-cover error. With $D_T=\min\{1,(N/c_*-1)/T\}$ and $e_T=(\kappa-c_*)D_T$, its realized risk is at most
$q_T+LD_T+Le_T/(N-\kappa+e_T)$. This difference tends to zero, while every dual value lower-bounds every feasible full-history strategy. Choose $T$ to make the difference at most $\eta$. The prior completion proof counts the finite prefix, fallback and initial-mixture roles. For variable dictionaries use the preceding decision-net or algebraic optimization, not convexification of their union.
\end{proof}


\begin{theorem}[Simultaneous expected resources]\label{thm:multi-resource}
The compact-model certificate and completion extend to finitely many additive nonnegative resource costs. Write $C_{\theta,T}^{(\ell)}$ for the prefix cost of resource $\ell=0,\ldots,r$, assume conditional lower and upper bounds $0\le c_\ell I\preccurlyeq G_{\theta,a}^{(\ell)}\preccurlyeq c_\ell^*I$, and require $c_0>0$. Fix a budget vector $\mathbf N$ and one common finite fallback with uniform resource bounds $\kappa_\ell<N_\ell$ simultaneously. Each finite constraint is
\[
 C_{\theta,T}^{(\ell)}+c_\ell S_{\theta,T}
                         \preccurlyeq N_\ell I.
\]
Every finite optimal plan admits an exact certificate with risk mass one and cost atoms indexed by a resource, a model and a pure preparation. If $\lambda_\ell$ is the total multiplier mass for resource $\ell$, then
\[
 \sum_{\ell=0}^r\lambda_\ell(N_\ell-\kappa_\ell)\le L,
 \qquad \text{total positive cost and risk atoms}\le k_T+1.
\]
Thus the support bound depends on the affine plan dimension, not directly on the number of resource constraints. The common Bellman recursion adds all weighted immediate resource costs; its reserve sums $c_\ell$ times the weighted open mass of the cost atoms for resource $\ell$.

Let $q_T(\mathbf N)$ denote the finite lower and put
\[
 D_T=\min\{1,(N_0/c_0-1)/T\},\quad
 e_{\ell,T}=(\kappa_\ell-c_\ell)D_T,\quad
 p_T=\min_{0\le\ell\le r}
 \frac{N_\ell-\kappa_\ell}{N_\ell-\kappa_\ell+e_{\ell,T}}.
\]
There is a common proper finite-label observer satisfying every original budget and having worst loss at most
\[
 q_T(\mathbf N)+LD_T+L(1-p_T).
\]
The difference tends to zero. At any feasible finite vector, primal attainment and monotone finite lower convergence also hold, without asserting a finite boundary implementation. The costs may represent separately charged calls, time, energy or communication, when they obey these additive hypotheses; a hard program-size, scratch-space or dictionary-cardinality constraint is not such a cost by declaration.
\end{theorem}
\begin{proof}
For the finite problem set $c_\ell(x)=\max_{\theta,u}\langle u,(C_{\theta,T}^{(\ell)}+c_\ell S_{\theta,T})u\rangle$ and $v(x)=\max_{\theta,u}\langle u,V_{\theta,T}u\rangle$. The upward closed set of achievable vectors $(c_0(x)-N_0,\ldots,c_r(x)-N_r,v(x))$ with coordinatewise slack is convex and closed, since $x$ lies in a compact polytope. A supporting hyperplane at $(0,\ldots,0,q_T)$ has nonnegative coefficients. Its risk coefficient is positive: if it were zero, the jointly strictly feasible $x_0$ would give a strictly smaller pairing. Normalize that coefficient to one. This yields multipliers $\lambda_\ell\ge0$, minimization of $v+\sum_\ell\lambda_\ell(c_\ell-N_\ell)$ and complementary slackness.

Repeat the active-gradient and normal-cone argument with one cost-gradient hull for each positive multiplier. All cost generators still have appended coordinate zero, and risk generators have appended coordinate one. The conic reduction in dimension $k_T+1$ therefore proves the same total support bound. Evaluating the resulting affine minimum at the jointly strict fallback prefix gives the displayed weighted margin bound. At each node the immediate scalar cost is simply the sum over resource-indexed atoms, which proves the recursion claim.

Resource 0 bounds survival by $D_T$. Completing an optimizer by the one common fallback costs at most $N_\ell+e_{\ell,T}$ in resource $\ell$ and adds at most $LD_T$ risk. Choose this completed plan with probability $p_T$ and the fallback otherwise. For every $\ell$ its cost is at most
$p_T(N_\ell+e_{\ell,T})+(1-p_T)\kappa_\ell\le N_\ell$.
The risk bound follows exactly as for one resource. Strict positive margins make $p_T\to1$. Common-tree and fallback labels are counted as before. On a general feasible vector, intersect all prefix reserve constraints; the single coercive resource 0 forces properness and uniform bounded-loss tails. The compactness and finite-intersection proof of Theorem~\ref{thm:boundary-attainment} then applies with all the additional closed constraints. Separate fallbacks safe for separate resources would not justify this proof: joint feasibility of one common fallback is essential.
\end{proof}

\subsection{Distinct witnesses and the price of hard output width}
\begin{example}[A shared adversarial prior can give a strict dual gap]\label{ex:distinct-dual-witnesses}
There are two commands $A,B$, one productive report, one decoder, and two scalar models. Costs are charged but are not extra observations; the visible action--report history has the same law in both models. Model 1 has costs $(1,2)$ and losses $(0,0)$; model 2 has costs $(1,1)$ and losses $(1,0)$. Commitment may occur only after a call. At budget $N=3/2$, writing $x$ for the probability that the final command is $B$, every plan costs at least $1+x$ in model 1 and has risk $1-x$ in model 2. One-call mixing at $x=1/2$ attains the exact value $1/2$. The dual certificate has one unit cost atom at model 1 and one unit risk atom at model 2.

If one incorrectly forces both atom lists to be proportional to the same probability $(t,1-t)$, the one-call dual becomes
\[
 \sup_{0\le t\le1,\lambda\ge0}\min_{0\le x\le1}
 \{(1-t)(1-x)+\lambda(tx-1/2)\}
       =\tfrac32-\sqrt2<\tfrac12.
\]
For $t\le1/2$ its best value is zero; for $t>1/2$ the best multiplier is $(1-t)/t$, giving $3/2-t-1/(2t)$, maximized at $t=1/\sqrt2$. All quantities come from real cost and loss tables. A single least-favorable prior for one scalarized task is not a substitute for independently imposed robust cost and risk constraints.
\end{example}

\begin{proposition}[An exact hard-width convexification gap]\label{prop:width-dual-gap}
There is a one-model, one-dimensional physical experiment with two commands, compact decision set $D=\{0,1\}$ and bounded nonnegative loss, such that for $1\le N\le2$,
\[
 R_1(N)=\min\{1/2,2-N\},\qquad R_2(N)=(2-N)/2.
\]
Thus $R_1$ is attained and right-continuous but not convex. Its scalar Lagrange relaxation over all width-one programs equals $R_2$; at $N=3/2$ it has a gap $1/4$. Randomizing two different one-value dictionaries pays for a two-value dictionary, not for an uncharged initial seed.
\end{proposition}
\begin{proof}
Both commands always give one productive report, with costs $1,2$. Decision 0 has loss $1/2$ after either command. Decision 1 has loss $1$ after the cheap command and zero after the costly one. These are permissible prospective audit losses: the future audit can be the deterministic command symbol with the displayed loss table. This is a bounded general-loss example, not a claim about a squared-loss realization.

For the sole decoder 0 every proper program has loss $1/2$. For the sole decoder 1, if its final command is costly with probability $x$, expected total cost is at least $1+x$ and risk is $1-x$. One-call mixing realizes all these pairs, so optimization gives $R_1$. With both values, the smallest loss is $1/2$ after a cheap last command and zero after a costly one. The same cost lower and one-call realization give $R_2$.

Over width-one programs the penalized infimum is
$\min\{1/2+\lambda(1-N),\lambda(2-N)\}$: the cheap decoder-1 endpoint is dominated. Maximizing over $\lambda\ge0$ gives $(2-N)/2$, with equality of the two lines at $\lambda=1/2$. At $N=3/2$ the primal width-one value is $1/2$ and the relaxation is $1/4$. In the randomized implementation attaining the latter, a cheap call uses decoder 0 and a costly call uses decoder 1, so two distinct final values are genuinely required.
\end{proof}

The preceding results supply an exact variational and finite-witness completion of the primitive risk route. They do not establish that every causal family admits the geometric reference certificates, a common finite unknown-model predictive quotient, or optimal computational resources. Boundary primal attainment is compatible with the failure of uniform expected cost-tail convergence and with lack of a finite boundary implementation. Those independent questions are not resolved by a success certificate for a finite optimization problem.
